\section{Construction of the implicit function}\label{sec:construction}

\subsection{Finite subdivision replaces the mean value theorem}

\begin{lemma}[Segment estimates]\label{lem:segment}
The hypotheses imply
\begin{align}
 \norm{H(x,y)-H(x',y)}&\le M\norm{x-x'},\label{eq:xseg}\\
 \norm{H(x,y)-H(x,y')}&\le\tfrac32\norm{y-y'},\label{eq:yseg}\\
 \norm{T_x(y)-T_x(y')}&\le\tfrac12\norm{y-y'},
 \qquad T_x(y)=y-H(x,y).\label{eq:tcontract}
\end{align}
Consequently \(H\) is Lipschitz on \(\Omega\) with constant
\(K=M+3/2\) for the product norm.
\end{lemma}

\begin{proof}
Consider first rational endpoints \(y,y'\). Fix a precision \(p\)
and put \(D=\norm{y'-y}\). Choose the explicit positive integer
\[
 L_p=\max\bigl(1,\lceil D\,2^{\mu(p)}\rceil\bigr),
 \qquad y_i=y+\frac{i}{L_p}(y'-y).
\]
Convexity keeps every point in \(Y\). Every increment has norm
\(D/L_p\le2^{-\mu(p)}\). Write the Taylor formula at \(y_i\):
\[
 T_x(y_{i+1})-T_x(y_i)
  =(\Id-V(x,y_i))(y_{i+1}-y_i)-R_i,
 \qquad \norm{R_i}\le2^{-p}D/L_p.
\]
Summing the \(L_p\) inequalities gives
\[
 \norm{T_x(y')-T_x(y)}
 \le(\tfrac12+2^{-p})D.
\]
The extra term can be made smaller than any prescribed rational
observation error, so \eqref{eq:tcontract} follows. Using
\(\norm V\le1+\norm{\Id-V}\le3/2\) in the same finite subdivision
gives \eqref{eq:yseg}. Subdivision in the \(x\)-coordinates gives
\eqref{eq:xseg}.

Here the summation is genuinely finite arithmetic. At a fixed
auxiliary observation precision \(t\), replace every value and
matrix by its rational evaluator. The accumulated additional error
is bounded, for example, by
\[
 2^{-t}(2L_p+2D)
\]
in the vertical \(T\) estimate. Induct on the finite rational
partial sum with \(p,t,L_p\) fixed, and then remove this explicitly
bounded error. Section~\ref{sec:arith} supplies the rational
squeezing argument. Thus this proof does not use induction on
the assertion that a real estimate holds at every precision.

The estimates define the continuous extension from rational points
to named points. Applying the rational bounds to sufficiently fine,
clamped approximations to the endpoints proves the same bounds
for those points. Finally, change \(x\) and then \(y\); the sum of
the first two bounds is at most
\((M+3/2)\norm{(x,y)-(x',y')}\).
\end{proof}

\subsection{The invariant box and the actual finite program}

The center estimate and \eqref{eq:xseg} give
\[
 \norm{H(x,b)}
 \le\norm{H(a,b)}+M\norm{x-a}\le r/2.
\]
For \(y\in Y\), the contraction estimate now yields
\begin{equation}\label{eq:invariant}
 \norm{T_x(y)-b}
 \le\norm{T_x(y)-T_x(b)}+\norm{H(x,b)}
 \le\tfrac12r+\tfrac12r=r.
\end{equation}
Thus the exact map \(T_x\) preserves \(Y\).

Let \(P_Y\) clamp each coordinate to the interval
\([b_i-r,b_i+r]\). For rational vectors it is a primitive recursive
map, and
\begin{equation}\label{eq:projection}
 P_Y(v)\in Y,\quad P_Y(y)=y\ (y\in Y),\quad
 \norm{P_Y(v)-P_Y(w)}\le\norm{v-w}.
\end{equation}
The scalar proof checks the order of two rational numbers and the
two endpoints; summing no coordinates is needed for the maximum
norm.

Define the dyadic floor map coordinatewise by
\[
 R_p(v)_i=2^{-p}\lfloor 2^p v_i\rfloor.
\]
It satisfies \(\norm{R_p(v)-v}\le2^{-p}\), also for negative
coordinates. Fix
\begin{equation}\label{eq:schedule}
 s=\lgup(2r),\qquad
 N(k)=k+s+4,\qquad
 \varepsilon_k=2^{-(k+4)}.
\end{equation}
For rational \(x\in X\), run exactly \(N(k)\) steps:
\begin{align}
 y_0^{x,k}&=b,\nonumber\\
 y_{j+1}^{x,k}
 &=P_Y\!\left(R_{k+5}
       \left(y_j^{x,k}-\Hhat((x,y_j^{x,k}),k+5)\right)\right),
       \quad j<N(k),\label{eq:algorithm}\\
 \Ghat(x,k)&=y_{N(k)}^{x,k}.\nonumber
\end{align}
This is primitive recursion in \(j\), followed by substitution
\(j=N(k)\). It never searches for a small residual.
The fixed precision and rounding also prevent gratuitous denominator
growth in the polynomial examples.

By the evaluation error, the rounding error, \eqref{eq:projection},
and \eqref{eq:invariant},
\begin{equation}\label{eq:step-error}
 y_j^{x,k}\in Y,\qquad
 \norm{y_{j+1}^{x,k}-T_x(y_j^{x,k})}
 \le2^{-(k+5)}+2^{-(k+5)}=\varepsilon_k.
\end{equation}
Projection does not change the exact image \(T_x(y_j)\), since
that image is already in the box. This fact is what makes
clipping a legitimate way to control evaluation and rounding
errors.

\subsection{A finite inexact contraction inequality}

\begin{lemma}[Inexact finite trace]\label{lem:recurrence}
If nonnegative rational numbers \(d_0,\ldots,d_N\) satisfy
\(d_{j+1}\le d_j/2+\eta\), then
\begin{equation}\label{eq:finite-bound}
 d_j\le2^{-j}d_0+2\eta(1-2^{-j})\le2^{-j}d_0+2\eta
 \quad(0\le j\le N).
\end{equation}
\end{lemma}

\begin{proof}
Use the quantifier-free invariant
\begin{equation}\label{eq:scaled-inv}
 2^j d_j+2\eta\le d_0+2^{j+1}\eta.
\end{equation}
For \(j=0\) this is equality. Multiplying the recurrence by
\(2^{j+1}\) and adding \(2\eta\), the left side at \(j+1\)
is at most
\[
 2^j d_j+2^{j+1}\eta+2\eta
 \le d_0+2^{j+2}\eta.
\]
All rational denominators may be cleared by a fixed positive
common denominator for the target finite trace. The induction
formula is then a natural-number inequality. Division by the
positive rational \(2^j\) gives \eqref{eq:finite-bound}.
\end{proof}

\subsection{Coherence by aligning the ends of two finite traces}

Let \(k\le l\). Compare the coarse run of length \(N(k)\)
with the last \(N(k)\) steps of the finer run of length \(N(l)\).
The initial points of these two pieces are in \(Y\), so their
distance is at most \(2r\). If the parameter values are \(x,x'\),
respectively, \eqref{eq:xseg}, \eqref{eq:tcontract}, and
\eqref{eq:step-error} give
\[
 d_{j+1}\le\tfrac12d_j+
 M\norm{x-x'}+\varepsilon_k+\varepsilon_l.
\]
Lemma~\ref{lem:recurrence} therefore proves
\begin{align}
 \norm{\Ghat(x,k)-\Ghat(x',l)}
 &\le2r\,2^{-N(k)}
       +2M\norm{x-x'}+2\varepsilon_k+2\varepsilon_l\nonumber\\
 &\le2M\norm{x-x'}+3\varepsilon_k+2\varepsilon_l.
 \label{eq:pairbound}
\end{align}
The last inequality uses \(2r\le2^s\).
Every distance in this comparison is between rational outputs of
finite programs.

Taking \(x=x'\), we obtain
\[
 \norm{\Ghat(x,k)-\Ghat(x,l)}
 \le5\varepsilon_k=\frac5{16}2^{-k}\quad(k\le l).
\]
Hence \(k\mapsto\Ghat(x,k)\) is a regular name in the sense of
\eqref{eq:regular}. \emph{This name is the definition of \(g(x)\).}
There is no appeal to an axiom asserting the existence of limits
of arbitrary Cauchy sequences. Comparing with an arbitrarily
finer stage in \eqref{eq:pairbound} gives the useful representation
error
\begin{equation}\label{eq:representation}
 \norm{\Ghat(x,k)-g(x)}\le3\varepsilon_k.
\end{equation}
As always, this means that the corresponding rational inequalities
hold at all requested finite observation precisions.

\subsection{Zero residual without an ideal infinite orbit}

For one fixed run, let
\(e_j=\norm{y_{j+1}^{x,k}-y_j^{x,k}}\).
We have \(e_0\le2r\), and, whenever the next term is present,
\[
 e_{j+1}\le\tfrac12e_j+2\varepsilon_k.
\]
Applying Lemma~\ref{lem:recurrence} to these finitely many
increments and using \(N=N(k)\ge1\) gives
\begin{align}
 \norm{H(x,\Ghat(x,k))}
 &=\norm{y_N-T_x(y_N)}\nonumber\\
 &\le\varepsilon_k+\tfrac12\norm{y_N-y_{N-1}}\nonumber\\
 &\le2r\,2^{-N}+3\varepsilon_k
 \le4\varepsilon_k=2^{-(k+2)}.
 \label{eq:residual}
\end{align}
By \eqref{eq:yseg} and \eqref{eq:representation},
\[
 \norm{H(x,g(x))}
 \le \tfrac32(3\varepsilon_k)+4\varepsilon_k
 =\tfrac{17}{2}\varepsilon_k
 \qquad\hbox{for every }k.
\]
For any observation error, an explicitly larger \(k\) makes the
right side smaller than that error. This proves \(H(x,g(x))=0\).
The box membership of the name follows from membership of every
term in the closed rational box.

\subsection{Uniqueness, residual bounds, and parameter dependence}

For \(y\in Y\), use \(T_x(g(x))=g(x)\) to obtain
\[
 \norm{y-g(x)}
 \le \norm{H(x,y)}+\norm{T_x(y)-T_x(g(x))}
 \le \norm{H(x,y)}+\tfrac12\norm{y-g(x)}.
\]
Rearrangement proves \eqref{eq:apost}. If \(H(x,y)=0\), the
distance is zero at every precision, proving uniqueness. At
\(y=b\) the same inequality and \eqref{eq:xseg} prove
\eqref{eq:location}.

Similarly,
\[
 \norm{g(x)-g(x')}
 \le\tfrac12\norm{g(x)-g(x')}+M\norm{x-x'},
\]
which proves \eqref{eq:glip}. Alternatively, it follows directly
by eliminating the explicit errors in \eqref{eq:pairbound}.
This latter proof uses only the two finite rational traces.

More generally, if \(\widetilde H\) has the same contraction and
box properties and its solution is \(\widetilde g\), then
\[
 \norm{g(x)-\widetilde g(x)}
 \le 2\norm{H(x,\widetilde g(x))-\widetilde H(x,\widetilde g(x))}.
 \tag{Perturbation}
\]
Thus a uniform normalized equation error \(\delta\) changes the
branch by at most \(2\delta\). This is another form of the
quantitative uniqueness behind the construction.

\subsection{Evaluation at an arbitrary primitive recursive parameter}

Suppose \(x(k)\) is a regular name of a point in \(X\).
Let \(C=\max(1,2M)\), \(c=\lgup(C)\), and set
\[
 q_k=P_X(x(k+c+2)),\qquad
 G_x(k)=\Ghat(q_k,k+2).
\]
The extension error is at most
\[
 2M\,2^{-(k+c+2)}+3\,2^{-(k+6)}
 \le \frac{19}{64}2^{-k}<2^{-k}.
\]
Equivalently, the pair estimate \eqref{eq:pairbound} directly
proves that \(G_x\) is a regular name independent of the chosen
name of \(x\). The Lipschitz bounds for \(H\) pass the zero
equation to this name. This construction uses only composition
with the fixed input name and primitive recursion already used
in \eqref{eq:algorithm}.
